\begin{lemma}[Upper enumeration in large induced subgraphs]
For every $\varepsilon>0$ and $0<d\le1$ there is $n_0$ such that the following holds.  Let $G[U,W]$ be an $(\varepsilon,d)$-regular pair with $|U|=|W|=n\ge n_0$.  If $S\subseteq U$, $T\subseteq W$, $S\cap T=\varnothing$, $|S|=|T|=y$, and $y\ge\sqrt\varepsilon n$, then the number of perfect matchings in $G[S,T]$ is at most
\[
 (d+\varepsilon+3\sqrt\varepsilon)^y y!.
\]
Here a missing edge can only reduce the number, so the assertion also applies to any balanced subgraph induced by prescribed subsets of the two parts.
\end{lemma}
\begin{proof}
Write $c=d+\varepsilon+3\kappa$ and let $m$ denote the number of perfect
matchings. If $c\ge1$, take $n_0=0$. The hypotheses $|S|=|T|=y$ and
\noderef{PerfectMatchingBijectionBound} give $m\le y!$, and
$c^y\ge1$ gives $m\le c^y y!$. This argument also covers $y=0$ and
infinite parameter values; the zeroth power is one.

Assume now $c<1$. All parameters are finite: $d,\varepsilon$ are bounded
by $c$, and $3\kappa\le c$ bounds $\kappa$. The identity
$\kappa^2=\varepsilon>0$ implies $\kappa>0$, and $c<1$ implies
$\kappa<1/3$. Consequently $\varepsilon=\kappa^2<\kappa$.
Conversion to real numbers preserves these finite nonnegative
parameters, their sums, products, inequalities, and powers. In the rest
of this branch use real notation for them, and put $D=d+\varepsilon>0$.

Apply \noderef{FactorialRootExceptionalProduct} to $D,\kappa$, obtaining
a positive integer $y_0$ depending only on those parameters. Choose an
integer $n_0>y_0/\kappa+1$. For $n\ge n_0$ and $y\ge\kappa n$ we have
$n>0$ and $y>y_0$. In particular $y>0$, so the empty induced graph
cannot occur in this branch.

Index $S$ and $T$ by $\{0,\ldots,y-1\}$, and write
$r_i=|N_G(e_S^{-1}(i))\cap T|$. Each $r_i\le y$, since its neighbor set
is a subset of $T$. To bound the exceptional indices, note that
$\varepsilon=\kappa^2<\kappa<1/3$, so $\varepsilon\le1$.
Since $n>0$, we have
$|T|=y\ge\kappa n>\varepsilon n=\varepsilon|W|$.
Apply \noderef{RegularPairUpperTail} with left part $U$, right part $W$,
and tested subset $T$, using the given regularity and $d\le1$.
It bounds the full $U$-tail
$B=\{v\in U:|N_G(v)\cap T|>(d+\varepsilon)y\}$ by
$|B|\le\varepsilon|U|=\varepsilon n$.
Let $A=\{v\in S:|N_G(v)\cap T|>(D+\kappa)y\}$.
Since $S\subseteq U$ and $(D+\kappa)y\ge Dy$, we have $A\subseteq B$.
The indexing bijection $e_S^{-1}$ restricts to a bijection from
$\{i:r_i>(D+\kappa)y\}$ onto $A$, so its cardinality $k$ equals $|A|$.
All thresholds are finite: $d+\varepsilon$ is finite, and all set
cardinalities are natural numbers. Their finite extended-real products
convert to the displayed real products, and strict comparison is
preserved by this conversion. Thus the real inequality defining $A$
implies the extended-real inequality defining $B$; converting the
extended-real tail bound back to reals gives
\[
 k=|A|\le |B|\le\varepsilon n
   =\kappa^2 n\le\kappa y,
\]
where the last step multiplies $\kappa n\le y$ by $\kappa>0$.

For the matching-product bound use \noderef{BregmanMatchingDegreeProduct}.
Its natural-number matrix has entry one exactly on the edges between
the indexed parts. Its zero-one hypothesis holds entrywise, its row
supports have cardinalities $r_i$, and its permanent is $m$ after the
natural-to-real cast, by \noderef{AdjacencyPermanentMatchingCount}.
That adapter applies the two branches of the checked imported theorem
\texttt{EntropyMethodsFormalization.BregmanMinc}. If a row has degree
zero, it gives $m=0$, and the assertion follows from nonnegativity.
Otherwise it gives
\[
 m\le\prod_{i=0}^{y-1}(r_i!)^{1/r_i}.
\]
All degrees are now positive and at most $y$, the exceptional count is
at most $\kappa y$, and $y\ge y_0$. Hence
\noderef{FactorialRootExceptionalProduct} gives
\[
 \prod_i(r_i!)^{1/r_i}\le(D+3\kappa)^y y!.
\]
In that estimate the increasing factorial roots bound exceptional rows
by $(y!)^{1/y}$; ordinary and low-degree rows are bounded by
$(y!)^{1/y}xe^{\delta_y}$, where
$x=D+2\kappa$ and $\delta_y=(\log y+1)/(\kappa y)$.
The positive margin
$\log((x+\kappa)/x)-\kappa\log(1/x)$ absorbs both this factorial loss
and the at most $\kappa y$ exceptional rows. Thus the regularity and
factorial losses, not just the imported permanent inequality, are
accounted for.

Combining the inequalities proves the real bound $m\le c^y y!$.
Apply the order-preserving map from nonnegative real numbers to
extended nonnegative real numbers. It preserves the products and
powers here and sends the natural-number count and factorial to their
usual casts. Since $c$ is finite it recovers precisely the original
coefficient. This yields the stated extended-real inequality. The
choice of $n_0$ involves only $D,\kappa$ and hence only the stipulated
parameters, and is uniform over the graph, the parts, and the subsets.
\end{proof}
